\section{Introduction}\label{sec:intro}
Deep learning~\cite{bib:lecun2015deep} (DL) models are increasingly deployed in safety- and mission-critical domains such as autonomous driving, health care, and energy, warranting research to ensure their quality in terms of accuracy, robustness, or efficiency~\cite{bib:sarker2021deep}.
Among existing quality assurance techniques for DL models, testing remains one of the most widely used approaches in practice~\cite{bib:zhang2022machine,bib:hu2024test}.

However, good performance on a test dataset does not guarantee model accuracy, robustness, or efficiency on novel inputs. Therefore, there is a need to systematically assess the quality of test datasets themselves to understand how thoroughly they evaluate the model.
Additionally, by quantifying test dataset quality, one can compare datasets to each other and pick one dataset over another in a given scenario.
For example, practitioners often face large pools of candidate test inputs but limited labeling budgets and computational resources. As a result, they need methods to identify a small subset of inputs that can adequately evaluate the model, allowing them to focus their effort on that subset.
Consequently, over the past decade, researchers have proposed numerous DNN test adequacy metrics to assess test-set quality and support applications such as test dataset prioritization and evaluation.
These metrics include, but are not limited to, neuron coverage~\cite{bib:pei2017deepxplore} and its variations~\cite{bib:ma2018deepgauge}, surprise adequacy~\cite{bib:kim2018guiding}, input distribution coverage~\cite{bib:dola2023input}, decision-boundary coverage~\cite{bib:liu2022deepboundary}, and source- and model-based mutation score~\cite{bib:hu2019deepmutation++,bib:humbatova2021deepcrime}.

Despite the large number of adequacy metrics, existing implementations remain highly fragmented. When the implementation is available, different approaches rely on substantially different preprocessing pipelines, execution workflows, and setup requirements.
For example, some metrics based on neuron coverage, like KMNC~\cite{bib:ma2018deepgauge}, require activation profiling over training data, surprise adequacy relies on activation trace extraction and density estimation, input distribution coverage depends on separately trained encoder models, and mutation testing typically requires generating and executing a large collection of mutants.
As a result, reproducing, comparing, and deploying these adequacy metrics in practice often requires substantial effort, either to re-implement unavailable metrics or to configure and run available ones.
Case in point, recent interviews with industrial DNN practitioners reveal that existing testing metrics and tools are difficult to use in practice due to heterogeneous workflows, high adoption barriers, limited tool integration, and the absence of publicly available implementations for some techniques~\cite{bib:you2025navigating}.

To address this problem, we present \tool, an extensible framework for running DL test adequacy metrics.
\tool provides a unified command-line interface for executing representative adequacy metrics, including neuron coverage and its variants, surprise adequacy, input distribution coverage, decision boundary coverage, and source-/model-level mutation score.
The framework is implemented in Python and supports Keras-based DL models.
Its unified architecture implements each adequacy metric as an independent module with consistent interfaces for preprocessing and scoring while preserving metric-specific execution logic, allowing users to easily integrate new metrics into the framework.
To improve efficiency under repeated evaluations, the framework also supports reusable caching of intermediate artifacts generated during preprocessing, which reduces redundant computation when metrics are repeatedly executed.
% To support practical usage, the framework provides command-line interfaces to evaluate test subsets within a consistent workflow.
In this paper, we present the engineering details of \tool and show how it can be used.

\tool is publicly available at \url{https://zenodo.org/records/21682100} and a demo video can be found at \url{https://aub.ie/ADEPT_video}.
% We demonstrate its usefulness through representative adequacy metrics and evaluation workflows drawn from practical DL testing scenarios.

% Our contributions are summarized as follows:

% \begin{itemize}

% \item We present \tool, a one-click, unified framework for DNN test dataset adequacy assessment that simplifies the execution and comparison of representative adequacy metrics through a consistent evaluation workflow.

% \item We integrate representative adequacy techniques with substantially different preprocessing requirements into a common execution framework, enabling consistent evaluation and comparison across neuron-coverage-based, surprise-adequacy-based, latent-distribution-based, decision-boundary-based, and mutation-score-based approaches.

% \item We provide standardized interfaces for metric execution, result collection, and runtime reporting, reducing engineering effort and improving the reproducibility and usability of DNN test dataset adequacy assessment.

% \end{itemize}